%% sections/results.tex
\section{Results}
\label{sec:results}

Every comparison holds the market structure, endowment draw, and seed
fixed, varying only the paradigm or plan.  We report five bubble
statistics following \citet{stockl2010bubble}: Haessel~$R^2$
\citep{haessel1978measuring}, normalised absolute deviation~$\ND$,
amplitude~$A$, turnover~$\TO$, and allocative efficiency~$\AllocEff$.
Formal definitions are in Appendix~\ref{app:instruments}.

\subsection{Strict-DLM replication}

\paragraph{Bubble attenuation}
In a typical session, prices overshoot $\FV_t$ by $20$--$40\%$ in
round~1, contract to $10$--$20\%$ in round~2, and track the
fundamental within $5\%$ by round~3.  Let
$\Delta_r=\max_t|\bar P_t{-}\FV_t|$ denote the maximum absolute
price deviation in round~$r$.  Across seeds,
$\Delta_1>\Delta_2>\Delta_3$ consistently.  The two-branch decision
rule of equation~\eqref{eq:dlm-gate}---momentum-chasing when
$k_i{=}0$, fundamental-anchoring when $k_i{\ge}1$---is sufficient to
reproduce the DLM result.  No cognitive model of learning is required;
the procedural switch alone generates the attenuation.

\paragraph{Round-4 treatments}
Table~\ref{tab:dlm-batch} reports batch results from ten sessions.
T4 produces higher mean deviation, turnover, and trade count in
round~4 than T2, replicating DLM's central comparison: the fraction
of experienced agents determines whether the bubble rekindles.

\begin{table}[h]
\centering\small
\caption{Strict-DLM batch: five T2 and five T4 sessions}
\label{tab:dlm-batch}
\begin{tabular}{l cc c cc}
\toprule
& \multicolumn{2}{c}{T2 (R4-$\sfrac{2}{3}$)}
& & \multicolumn{2}{c}{T4 (R4-$\sfrac{1}{3}$)} \\
\cmidrule{2-3}\cmidrule{5-6}
& Mean & SD & & Mean & SD \\
\midrule
Mean deviation (R4) & 3.2 & 1.8 & & 8.7 & 3.5 \\
Turnover (R4)       & 0.42 & 0.11 & & 0.68 & 0.15 \\
Trades (R4)         & 14.6 & 3.2 & & 22.4 & 4.8 \\
Session payoff      & 4120 & 280 & & 3940 & 350 \\
\bottomrule
\end{tabular}
\end{table}

\subsection{Plan~I baseline}

Under Plan~I, utility agents update beliefs via
equation~\eqref{eq:plan1}.  Round-1 novices ($w{=}0.60$) weight peers
heavily; by round~3 ($w{=}0.90$) the prior dominates.  Amplitude
$A\in[0.20,0.40]$, $\ND\in[0.10,0.25]$, and $R^2$ rises to
$0.40$--$0.80$ across the session.  The bubble-crash signature
persists but is attenuated relative to pure $w{=}0.60$ runs,
confirming that the experience ramp anchors prices toward the
fundamental.

Risk-typed stratification is visible: risk-loving, optimistic agents
carry posteriors above $\FV_t$; risk-averse, pessimistic agents carry
posteriors below~$\FV_t$.  Allocative efficiency
$\AllocEff\in[0.70,0.95]$.

\subsection{Plans II and III}

Under Plan~II, the LLM's posteriors are heterogeneous across risk
types in the expected direction and produce slightly higher amplitude
and lower $R^2$ than Plan~I: the LLM posterior is less tightly
anchored to the prior.

Under Plan~III, matched-seed comparisons show posterior distributions
that overlap heavily with Plan~II.  Per-period mean prices track
within a few cents; $A$, $\ND$, and $R^2$ are indistinguishable.

Table~\ref{tab:mktq} summarises.

\begin{table}[h]
\centering\small
\caption{Market-quality statistics across paradigms and plans
(round-4 for Strict-DLM; session-wide for Plans~I--III)}
\label{tab:mktq}
\begin{tabular}{l cccc}
\toprule
& Strict-DLM & Plan~I & Plan~II & Plan~III \\
\midrule
$R^2$       & .80--.95 & .40--.80 & .30--.75 & .30--.75 \\
$\ND$       & .05--.15 & .10--.25 & .10--.30 & .10--.30 \\
$A$         & .10--.25 & .20--.40 & .25--.45 & .25--.45 \\
$\TO$       & .3--.7   & .5--1.0  & .5--1.0  & .5--1.0  \\
$\AllocEff$ & ---      & .70--.95 & .60--.90 & .60--.90 \\
\bottomrule
\end{tabular}
\end{table}

Three findings emerge.  First, strict-DLM converges most cleanly
because the experienced branch's horizon discount provides the
strongest fundamental anchor.  Second, Plans~II and~III introduce
dispersion through the LLM reply distribution without producing a
qualitatively different regime.  Third, the Plan~II/Plan~III gap is
within measurement noise: the language model's internal representation
of ``risk-averse'' already encodes concavity comparable to
$\sqrt{w/w_0}$.  Naming the formula adds redundant, not corrective,
information.

\subsection{Robustness}

\paragraph{Background agents}
Removing the fundamentalist raises amplitude by $30$--$50\%$;
removing the trend follower suppresses bubbles almost entirely.  The
background agents play their intended structural roles.

\paragraph{Risk shares}
Skewing the population toward risk-loving agents raises mean prices
and amplitude; skewing toward risk-averse agents lowers both and
accelerates convergence.

\paragraph{Plan~I weight schedule}
Raising $w_{\text{naive}}$ above 0.60 suppresses the round-1 bubble;
lowering $w_{\text{skeptical}}$ below 0.90 prevents late-round
convergence.  The default $(0.60,0.90)$ best matches DLM's
round-by-round attenuation.

\paragraph{EU filtering}
The expected-utility score~\eqref{eq:eu} compresses the belief
distribution before it reaches the order book: an agent with an
extreme posterior but low cash still holds.  This filtering is the
primary reason the belief-update channel has a muted impact on
market-level statistics.
